\section{Background}

The emergence of \textbf{multi-agent AI systems} as a practical deployment paradigm introduces a constellation of engineering challenges that were largely theoretical a decade ago. Chief among these is the problem of \textbf{agent lifecycle management} — how autonomous agents are spawned, delegated to, monitored, and terminated within a coordinated system — and the parallel problem of \textbf{accountability and provenance}, which asks: when an agent chain produces a decision or action, who is responsible, and how can the chain of custody be reconstructed with confidence? These two problems are deeply intertwined. A system that spawns sub-agents without maintaining an immutable record of the parent-child relationship cannot later answer questions about responsibility, intent, or data origin. This section surveys the relevant literature across five areas: agent lifecycle management in multi-agent systems (\S\ref{sec:lifecycle}), accountability and provenance in distributed AI (\S\ref{sec:provenance}), fixed versus mutable delegation chains (\S\ref{sec:delegation}), hierarchical task decomposition in AI (\S\ref{sec:htd}), and the BX3 Framework as the substrate on which immutable parent chains become architecturally feasible rather than merely policy-dependent (\S\ref{sec:bx3}).

% ---------------------------------------------------------------------
% 2.1 Agent Lifecycle Management
% ---------------------------------------------------------------------
\subsection{Agent Lifecycle Management in Multi-Agent Systems}
\label{sec:lifecycle}

Modern multi-agent architectures decompose complex objectives across hierarchies of specialized agents, where a parent orchestrator delegates sub-tasks to spawned agents with potentially different capabilities, contexts, and access scopes. Variants of this pattern appear in LangChain's agent executor~\cite{langchain2023}, CrewAI's role-based agents~\cite{crewai2024}, and the MAKER medical-diagnosis pattern — and have revealed that lifecycle management is far more than a deployment detail: it determines whether agent systems are debuggable, auditable, and safe to operate at scale.

The central lifecycle challenge is that agents are not static processes. They may be created dynamically, inherit partial context from a parent, execute asynchronously, and produce results that feed into subsequent agents. \textbf{AgentSpawn} demonstrates that adaptive multi-agent collaboration through dynamic spawning achieves approximately 34\% higher task-completion rates on long-horizon code generation benchmarks, compared to static agent pools~\cite{agentspawn2026}. The spawn mechanism design — whether blocking or non-blocking, whether context is deep-copied or referenced, whether the child is ephemeral or persistent — has downstream consequences for both reliability and accountability. AgentSpawn manages memory transfer at spawn time, transmitting only context relevant to the child's assigned sub-task and reducing token overhead by approximately 42\% relative to full-context transfers.

Practical production systems have converged on a layered architecture comprising: (1) a \textbf{coordinator layer} responsible for task decomposition and agent spawning; (2) a \textbf{tool layer} that exposes sub-agents as composable tools to parent agents; (3) an \textbf{execution layer} that handles asynchronous result aggregation; and (4) a \textbf{policy layer} that enforces permission inheritance, rate limiting, and checkpointing~\cite{rde2024}. Critical lifecycle operations — cascading termination, orphan detection, cost tracking, and checkpoint/resume — are identified as production-hardening requirements that most research systems neglect. Without explicit lifecycle primitives, spawned agents that encounter errors or timeouts become orphaned processes whose state is unrecoverable and whose actions are untraceable.

The \textbf{Agent Lifecycle Protocol (ALP)} defines a formal, interoperable standard covering birth (Genesis), authority splits (Fork), cross-environment movement (Migration), capability updates (Retraining), obligation handoffs (Succession), and retirement (Decommission)~\cite{wooldridge1995,stone2000,ferber1996,padgham2001}. ALP introduces a genealogical registry capturing both genetic lineage (model and architecture provenance) and epigenetic lineage (configuration, memory, and reputation state), enabling full traceability across agent generations. This is directly relevant to recursive spawning: each spawned child inherits a lineage record tracing back through all ancestor agents, creating an auditable family tree of every delegation event.

The POSIX-inspired \textbf{Quine} framework maps agent identity to OS processes (PID), interface to standard streams, and lifecycle to fork/exec/exit primitives~\cite{zhuge2024}. Strong kernel-level isolation between parent and child agents is achieved by leveraging mature OS scheduling and fault-containment mechanisms. Quine supports recursive self-spawning, where an agent can \texttt{exec} a fresh instance of itself for self-renewal without a bespoke runtime. The parent-child relationship is tracked via the process tree, and termination semantics are inherited from the kernel.

The \textbf{AGENTHIVE} framework treats delegation as a first-class, recursive primitive where any agent can spawn sub-agents without a fixed central orchestrator~\cite{agenthive2025}. Resulting topologies — trees, forests, star structures — are emergent rather than pre-configured, driven by task demands. Mutable, demand-driven models enable highly adaptive structures but introduce challenges for auditability and fixed-chain guarantees. This trade-off is resolved in the BX3 Framework by constraining AGENTHIVE-style recursive spawning within an immutable parent-substrate layer.

% ---------------------------------------------------------------------
% 2.2 Accountability and Provenance
% ---------------------------------------------------------------------
\subsection{Accountability and Provenance in Distributed AI Systems}
\label{sec:provenance}

The accountability problem in multi-agent systems is not merely governance philosophy — it is a technical challenge in distributed systems design. An auditability framework for LLM-based agent systems identifies five dimensions that determine whether a system can be held accountable post-deployment: action recoverability, lifecycle coverage, policy checkability, responsibility attribution, and evidence integrity~\cite{trusttrack2025}. Of these, the most operationally complex is \textbf{responsibility attribution} in multi-agent chains: when Agent C produces an output tracing back through Agent B to Agent A (spawned by a human), the system must attribute the full chain of decisions and actions to the appropriate principals.

A parallel concern is \textbf{provenance} — the tamper-evident record of what data, models, and governance controls influenced each decision. Provenance-aware AI systems propose structured audit trace collection, capturing not just inputs and outputs but the full context of agent reasoning and delegation decisions~\cite{aer2025}. The EU AI Act's record-keeping expectations for high-risk systems have accelerated interest in provenance architectures that can survive regulatory scrutiny. Prior work reports median enforcement costs of approximately 8.3\,ms for pre-execution mediation with tamper-evident logging, suggesting that strong accountability guarantees are achievable without prohibitive overhead~\cite{vap2026}.

\textbf{Agent Execution Records (AER)} formalize first-class provenance primitives for autonomous agents, proposing action-state-attribution tuples that record not just what an agent did but which parent agent authorized it, under which declared scope, and at what point in the chain~\cite{aer2025}. \textbf{Warrant Certificate Authorities (WCA)} extend this with auditable data provenance for tool-call chains, establishing that each tool invocation can be cryptographically attributed to a specific warrant issued by a parent agent~\cite{wca2025}. The \textbf{TIBET} protocol further proposes transaction/interaction-based evidence trails for AI-to-AI and AI-to-human interactions, creating a unified evidentiary chain across heterogeneous agent topologies~\cite{tibet2026}.

The \textbf{Human Delegation Provenance (HDP)} protocol is the work most directly relevant to the contribution of this paper. HDP proposes a lightweight token-based scheme that captures human authorization decisions and propagates them through the delegation chain as first-class cryptographic artifacts~\cite{hdp2026}. The protocol's core insight is that human authorization is not a single permission granted at chain initialization but a per-decision artifact that must travel with each delegation event. Every HDP token encodes scope, duration, and constraints, and cryptographic verification is performed at each chain traversal. However, HDP assumes a trustworthy delegation chain topology — it does not structurally enforce the integrity of the chain itself. A compromised or misbehaving agent can present a false parent identity before presenting tokens for verification. The combination of HDP-style per-decision provenance with a structurally immutable parent chain closes both gaps simultaneously.

The \textbf{MAIF} consortium proposes an artifact-centric agentic paradigm that enforces AI trust and provenance through tamper-evident artifact chains rather than runtime assertion~\cite{maif2025}. \textbf{Chain of Consciousness} introduces a cryptographic protocol for verifiable AI agent provenance using hash-chained provenance records that make retrospective modification computationally detectable~\cite{chainofconsciousness2026}. These approaches collectively establish that provenance is achievable in multi-agent systems but that its guarantees depend on the immutability of the underlying delegation topology — a property that existing frameworks treat as an assumption rather than an architectural primitive.

% ---------------------------------------------------------------------
% 2.3 Fixed vs. Mutable Delegation Chains
% ---------------------------------------------------------------------
\subsection{Fixed vs. Mutable Delegation Chains}
\label{sec:delegation}

The question of whether a delegation chain should be fixed at spawn time or allowed to mutate dynamically is a foundational design decision with consequences for both security and auditability. The dominant pattern in existing frameworks is \textbf{mutable delegation}, where parent agents may re-delegate, escalate, or transfer authority to other agents during execution. This provides flexibility but creates a significant provenance problem: if the delegation chain can change at runtime, the historical record of who authorized what becomes ambiguous, and post-hoc auditing requires reconstructing not just what happened but what the authorization state was at each decision point.

The \textbf{Intent-Preserving Delegation Protocol (IPDP)} addresses mutable chains through a non-LLM Delegation Authority Service that enforces pre-execution intent checking, at-execution scope verification, and post-execution evidence logging~\cite{sentinelagent2025}. SentinelAgent's Delegation Chain Calculus (DCC) defines seven properties for verifiable delegation in federal multi-agent AI systems, including \textit{forensic reconstructibility} (the chain can be completely reconstructed from logs) and \textit{authority narrowing} (each hop can only reduce, never expand, the authority of the previous hop).

A fixed delegation chain can be viewed as a constrained case of a mutable chain where the mutation log is empty. The advantage of the fixed model is simplicity and performance — no hop-resolution overhead. The advantage of the mutable model is flexibility and fault tolerance. The BX3 Framework resolves this tension by treating the delegation chain as immutable at the level of parent identity while allowing the child agent full autonomy over its own execution and further spawning. The practical cost of fixed chains is real but manageable: load rebalancing in a fixed-chain system requires spawning new nodes rather than redirecting existing ones; a failed agent cannot be transparently replaced by reassigning its children. In high-stakes environments — legal reasoning, medical diagnosis, financial compliance, critical infrastructure orchestration — temporary throughput reduction during a failure incident is far less costly than accountability loss.

Existing approaches to chain integrity management are uniformly policy-based. Access-control lists, role-based security models, and audit logs are all enforced by the runtime's execution logic~\cite{wooldridge1995}. The fundamental vulnerability of policy-based chain management is that it trusts the agent to enforce constraints against its own operational interests. An agent near completion of a valuable task has a strong incentive to suppress escalation, exceed its token budget, or silently absorb a failure rather than surface it to the human anchor and risk termination. The BX3 Framework addresses this by making the parent pointer a structural property enforced by the \fact{} layer — not a policy parameter that the agent can rewrite.

\textbf{Agent Contracts} formalize governance for autonomous AI agents through bilateral agreements between principals and agents, encoding obligations, scope, and revocation conditions as first-class syntactic objects~\cite{agentcontracts2025}. \textbf{Proof of Bottleneck (PoB)} introduces a cryptographic protocol where a designated Bottleneck agent must approve all cross-domain communications, creating a centralized audit point that makes cross-domain provenance verifiable without real-time external verification~\cite{vap2026}. The BX3 Framework complements these approaches with immutable parent references that make the delegation chain itself structurally tamper-evident — a property that PoB's bottleneck model requires but does not independently guarantee.

% ---------------------------------------------------------------------
% 2.4 Hierarchical Task Decomposition
% ---------------------------------------------------------------------
\subsection{Hierarchical Task Decomposition in AI}
\label{sec:htd}

Multi-agent systems do not operate in isolation — their spawning behavior is triggered by the need to decompose complex tasks into manageable, parallelizable sub-tasks. \textbf{Hierarchical Task Network (HTN)} planning, originally developed in classical AI, has experienced renewed interest as a structuring principle for LLM-based agents, enabling localized replanning, parallel sub-task execution, and modular failure recovery~\cite{htn1993,htn2000}. The field has converged on complementary techniques: task decomposition strategies that break goals into sub-tasks, hierarchical planning structures that organize sub-tasks into executable trees, execution patterns that govern plan-real-world feedback interaction, and replanning mechanisms that handle divergence between plans and reality.

A critical observation from recent work is that static decomposition is insufficient. The \textbf{ADaPT} framework introduces as-needed decomposition, where a separate planner module leverages an LLM to create high-level plans only when the executor agent cannot directly handle a sub-task — recursively decomposing complex sub-tasks on demand rather than pre-planning the entire hierarchy~\cite{recap2025}. \textbf{ReAcTree} extends this with a dynamic tree structure where each agent node retrieves goal-specific experiences from episodic memory and shares information via working memory, enabling hierarchical task planning that adapts to execution feedback in real time~\cite{reactree2025}. \textbf{DeepAgent} proposes a hierarchical task DAG-based autonomous framework where the graph expands tasks into finer sub-tasks as new information arrives, enabling real-time adaptation to changing task context~\cite{deepagent2025}. \textbf{HiPlan} provides hierarchical planning for LLM agents with adaptive global-local guidance, and \textbf{StructuredAgent} plans with AND/OR trees for long-horizon web tasks~\cite{hiplan2025,structuredagent2025}.

The common thread across this literature is that effective task decomposition is not purely a planning problem — it is deeply entangled with the lifecycle and accountability of the agents that execute the decomposed sub-tasks. A decomposition strategy that spawns agents without maintaining an immutable parent chain produces a tree of actions that cannot be audited to its root. Conversely, a system with immutable parent chains but primitive decomposition logic will spawn agents inefficiently. The synthesis of these concerns — hierarchical decomposition with immutable delegation provenance — is the open problem that the BX3 Framework's Recursive Spawning Protocol addresses.

% ---------------------------------------------------------------------
% 2.5 BX3 Framework as Substrate for Immutable Parent Chains
% ---------------------------------------------------------------------
\subsection{The BX3 Framework as Substrate for Immutable Parent Chains}
\label{sec:bx3}

The four research threads surveyed above converge in the BX3 Framework's architectural response to recursive spawning. BX3 — a universal functional-role architecture defining how agents are structured, communicate, and govern themselves across the Agentic platform and all Bxthre3 ventures — establishes an immutable parent-chain substrate that addresses the accountability, provenance, and delegation challenges identified in the preceding sections~\cite{bx3framework,bx3framework2026}.

BX3 defines three functional layers operating within every agent node: the \textbf{Purpose Layer} (carrying intent, judgment, and final human accountability), the \textbf{Bounds Engine} (carrying constrained execution within a defined operating range), and the \textbf{Fact Layer} (carrying deterministic enforcement and forensic auditability)~\cite{bx3framework2026}. Each layer is defined by the properties it must maintain rather than by the type of actor occupying it. This decomposition is not organizational — it is a physical constraint enforced at the interface between the reasoning system and the physical world.

The key architectural property for recursive spawning is that the parent reference is established \textbf{exactly once} at spawn time via a Fact Layer atomic write, and is thereafter immutable. The reference encodes the parent's canonical identity, the spawn timestamp, the declared scope of delegation, and a SHA-256 digest of the parent's state at the time of spawn. This design provides forensic traceability: any action taken by a child agent can be traced back through its full ancestor chain to the originating human principal or root orchestrator. No user, agent, or system call can alter the parent field after creation.

The \textbf{Sandbox Gate Protocol} enforces the boundary between parent and child execution contexts. The parent agent defines a constrained execution environment for the child — permissible tool calls, argument constraints, and scope limits — and the child cannot expand its own privileges beyond what the parent has declared. This mirrors the monotonic attenuation property of authorization attestation tokens and the non-widening delegation constraints of the Agent Passport System~\cite{sovereignagent}, extended with BX3-specific role semantics that map permissions to functional role definitions rather than raw capability flags.

The \textbf{Forensic Ledger Architecture} provides the tamper-evident, append-only audit trail required for accountable recursive spawning. Every spawn event, delegation declaration, tool invocation, and decision output is recorded in the forensic ledger with cryptographic linkage to the agent's immutable parent reference. Entries are chained by SHA-256 hash — each entry contains the hash of the preceding entry — creating a tamper-evident structure that makes retrospective modification of any entry computationally detectable.

The \textbf{Truth Gate} mechanism acts as the pre-execution policy gate for all agent actions within the spawn tree. Before any child agent commits to an action, the Truth Gate evaluates the proposed action against the declared delegation scope and the parent agent's safety constraints. DENIED receipts are logged in the forensic ledger with full provenance, enabling post-hoc detection of policy violations or scope creep.

The \textbf{Training Wheels Protocol} governs the graduated autonomy of spawned agents~\cite{sovereignagent}. In HITL-active Mode 1 (the default), every outbound action by any agent — including spawned children — is queued for human review before execution. Approvals and rejections are logged as behavioral patterns that inform future spawn decisions, enabling the parent agent to refine its delegation scope based on observed child reliability. This creates a closed feedback loop between provenance tracking, accountability, and adaptive spawning policy.

The integration of these mechanisms positions recursive spawning in BX3 as a governed, auditable, and verifiably accountable operation rather than an unchecked runtime expansion. Spawned child agents inherit immutable parent references, execute within role-scoped constraints enforced by the Sandbox Gate, produce tamper-evident ledger entries via Forensic Ledger Architecture, pass through Truth Gate evaluation before committing actions, and operate under Training Wheels oversight until earned autonomy is demonstrated. The result is a spawn tree that maintains the compositional flexibility of AGENTHIVE-style recursive spawning while delivering the accountability guarantees demanded by the provenance and auditability frameworks surveyed above.
